%% kapa: exam=1 round=28 year=2017 period=1 form=A subject=civil_law subject_ko=민법 seq=4 number=4
%% source: Q-Net 감정평가사 기출문제 (cst00309) artlSeq=5199603 file=extracted/1008/5199603/2017년 제28회 감정평가사 1차시험 문제지/1교시 A형.pdf
%% conversion: kordoc 4.24.51 -> tools/gen1.py (auto); review=proofread-sonnet; flags=none
\begin{question}{4}
부재와 실종에 관한 설명으로 옳지 않은 것은? (다툼이 있으면 판례에 따름)
\begin{choices}
  \item 부재자는 성질상 자연인에 한한다.
  \item 법원은 선임한 재산관리인에 대하여 부재자의 재산으로 상당한 보수를 지급할 수 있다.
  \item 외국에 장기 체류하는 자가 국내에 있는 재산을 관리하고 있으면 그는 부재자에 해당하지 않는다.
  \item 부재자에 대한 실종선고 이전에 법원이 선임한 부재자의 재산관리인이 선임결정 취소 전에 한 처분행위에 기하여 경료된 등기는 적법한 것으로 추정된다.
  \item 피상속인의 사망 후에 그의 아들에 대한 실종선고가 있었으나 피상속인의 사망 이전에 실종기간이 만료된 경우, 그 아들은 상속인이 될 수 있다.
\end{choices}
\end{question}
\answer{5}
